\chapter{Proposta da Solução} 
\label{cap:proposta} 

\section{Visão geral da solução} 
\label{sec:visao-geral-solucao} 
A solução proposta neste trabalho consiste em uma implementação funcional baseada em Modelos de Linguagem de Grande Porte, voltada à simplificação e adaptação de textos técnicos de saúde para diferentes perfis de pacientes. A proposta parte do problema apresentado na Introdução: conteúdos de saúde podem ser corretos do ponto de vista técnico, mas ainda assim difíceis de compreender quando são escritos em linguagem especializada ou pouco adequada ao nível de letramento em saúde do usuário final.

O sistema desenvolvido atua sobre conteúdos previamente definidos ou aprovados por profissionais ou equipes de saúde. Dessa forma, seu objetivo não é criar novas orientações clínicas, mas apoiar a transformação comunicacional de materiais já existentes. A solução busca adaptar a forma de apresentação desses conteúdos, tornando-os mais claros, diretos e acessíveis, sem alterar o significado do texto original. 

O funcionamento geral da solução é apresentado na Figura~\ref{fig:fluxo-conceitual-solucao}. O fluxo parte do conteúdo do plano de cuidado, do perfil comunicacional do paciente e de parâmetros definidos no momento do processamento. Esses elementos são combinados para orientar a geração das saídas pela LLM, mantendo a adaptação restrita ao conteúdo fornecido e ao objetivo de simplificação comunicacional.

\def\escalaFigConceitual{0.85}
\begin{figure}[htbp]
  \centering
  \resizebox{\escalaFigConceitual\textwidth}{!}{
  \begin{tikzpicture}[
    >=Latex,
    every node/.style={font=\small},
    startend/.style={
      rounded rectangle,
      rounded rectangle arc length=180,
      draw,
      align=center,
      minimum height=1.0cm,
      text width=4.8cm
    },
    process/.style={
      rectangle,
      draw,
      align=center,
      minimum height=1.0cm,
      text width=4.8cm
    },
    guardrail/.style={
      rectangle,
      draw,
      dashed,
      fill=gray!20,
      align=center,
      minimum height=1.0cm,
      text width=4.4cm
    },
    output/.style={
      rounded rectangle,
      rounded rectangle arc length=180,
      draw,
      align=center,
      minimum height=2.7cm,
      text width=6.6cm
    },
    arrow/.style={->, thick}
  ]

    %------------------------------------------------
    % LINHA SUPERIOR: entradas
    %------------------------------------------------
    \node[startend] (conteudo) at (0,0) {
      \textbf{Conteúdo do plano de cuidado}\\
      \textit{PDF ou texto livre}
    };

    \node[startend] (perfil) at (7.0,0) {
      \textbf{Perfil comunicacional}\\
      \textit{idade, escolaridade e área de formação}
    };

    \node[startend] (parametros) at (14.0,0) {
      \textbf{Parâmetros de processamento}\\
      \textit{nível de detalhamento e inclusão de imagens}
    };

    %------------------------------------------------
    % INSTRUÇÃO CENTRAL
    %------------------------------------------------
    \node[process] (instrucao) at (7.0,-3.2) {
      \textbf{Instrução para a LLM}\\
      \textit{texto original, perfil, parâmetros,}\\
      \textit{regras de simplificação e guardrails}
    };

    % Setas das entradas para a instrução
    \draw[arrow] (conteudo.south) |- (instrucao.west);
    \draw[arrow] (perfil.south) -- (instrucao.north);
    \draw[arrow] (parametros.south) |- (instrucao.east);

    %------------------------------------------------
    % RAMOS PARALELOS
    %------------------------------------------------
    \node[process] (geracaoTexto) at (3.4,-6.3) {
      \textbf{Geração do texto simplificado}
    };

    \node[process] (geracaoFAQ) at (10.6,-6.3) {
      \textbf{Geração de perguntas e respostas}
    };

    \coordinate (midBranch) at (7.0,-4.9);
    \draw[thick] (instrucao.south) -- (midBranch);
    \draw[arrow] (midBranch) -| (geracaoTexto.north);
    \draw[arrow] (midBranch) -| (geracaoFAQ.north);

    % Guardrails
    \node[guardrail] (guardTexto) at (3.4,-8.9) {
      \textbf{LLM com guardrails}\\
      \textit{preservação do significado}
    };

    \node[guardrail] (guardFAQ) at (10.6,-8.9) {
      \textbf{LLM com guardrails}\\
      \textit{restrição ao conteúdo original}
    };

    \draw[arrow] (geracaoTexto) -- (guardTexto);
    \draw[arrow] (geracaoFAQ) -- (guardFAQ);

    % Saídas imediatas
    \node[process] (texto) at (3.4,-11.4) {
      \textbf{Texto simplificado e adaptado}
    };

    \node[process] (faq) at (10.6,-11.4) {
      \textbf{Perguntas frequentes}
    };

    \draw[arrow] (guardTexto) -- (texto);
    \draw[arrow] (guardFAQ) -- (faq);

    %------------------------------------------------
    % SAÍDAS CONSOLIDADAS
    %------------------------------------------------
    \node[output] (saidas) at (7.0,-14.8) {
      \textbf{Saídas da solução}\\[0.12cm]
      \textit{texto simplificado e adaptado}\\
      \textit{perguntas e respostas}\\
      \textit{glossário de termos técnicos}\\
      \textit{Flesch-PT}\\
      \textit{chat contextual com guardrails}
    };

    \draw[arrow] (texto.south) -- ++(0,-0.8) -| ($(saidas.north)+(-1.6,0)$);
    \draw[arrow] (faq.south) -- ++(0,-0.8) -| ($(saidas.north)+(1.6,0)$);

  \end{tikzpicture}
  }
  \caption{Fluxo conceitual da solução proposta.}
  \label{fig:fluxo-conceitual-solucao}

  {\small Fonte: Elaboração própria.}
\end{figure}

A partir dessas entradas, a solução constrói instruções direcionadas à LLM, combinando o texto original, o perfil do paciente, os parâmetros de processamento, as regras de simplificação e os guardrails. O processamento ocorre em dois ramos paralelos. No primeiro, a LLM é utilizada para gerar uma versão simplificada e adaptada do texto. No segundo, a LLM é utilizada para gerar perguntas frequentes e respostas baseadas no conteúdo original. Em ambos os ramos, os guardrails delimitam a geração, restringindo as saídas ao conteúdo fornecido e orientando a preservação do significado original. 

A LLM, nesse processo, é tratada como mecanismo de reescrita e adaptação linguística, e não como fonte autônoma de conhecimento médico. A preservação da precisão médica é entendida como a manutenção do significado do texto original durante a simplificação e a geração dos recursos auxiliares. Assim, a solução não tem como finalidade diagnosticar, prescrever, recomendar novas condutas, alterar orientações médicas ou validar clinicamente o material de entrada. 

Como saída, a solução apresenta o texto simplificado e adaptado, recursos auxiliares de compreensão, como glossário e perguntas frequentes, chat contextual e um indicador automático de legibilidade textual. Neste trabalho, esse indicador é o Flesch-PT, utilizado para comparar o texto original e a versão simplificada do ponto de vista linguístico. Essa métrica apoia a avaliação da simplificação textual, mas não substitui a análise da preservação do conteúdo médico.

\section{Contexto de aplicação na Takere}
\label{sec:contexto-aplicacao-takere}

O fluxo apresentado na seção anterior corresponde ao funcionamento implementado neste TCC para demonstração e avaliação da funcionalidade proposta. A Takere é tomada como ambiente de referência por se tratar de uma plataforma no-code para o desenvolvimento de aplicações mHealth baseadas em planos de cuidado \cite{niemiec2022takere}. Assim, a plataforma não constitui o tema central do trabalho, mas o contexto no qual a solução poderia ser aplicada. 

Em aplicações baseadas em planos de cuidado, textos técnicos podem aparecer em diferentes elementos, como orientações, descrições de atividades, lembretes, perguntas e explicações apresentadas ao paciente. Embora esses conteúdos possam ser definidos por profissionais ou equipes de saúde, sua redação nem sempre é adequada ao nível de letramento em saúde do usuário final. Nesse cenário, a solução proposta atua como apoio à adaptação comunicacional desses textos, gerando versões mais claras e acessíveis sem modificar seu significado original. 

No contexto da Takere, a solução poderia ser incorporada futuramente ao fluxo de apresentação de conteúdos ao paciente, atuando sobre textos associados aos planos de cuidado. No entanto, a implementação desenvolvida neste trabalho não pressupõe integração plena com todos os módulos da plataforma; ela demonstra, de forma controlada, como a simplificação textual orientada por perfil e delimitada por guardrails pode ser aplicada a conteúdos de saúde previamente definidos.

\section{Entradas da solução} 
\label{sec:entradas-solucao} A solução recebe como entrada principal um conteúdo de saúde relacionado a um plano de cuidado. Esse conteúdo pode ser fornecido de duas formas: por meio de um arquivo PDF ou pela inserção direta de um texto livre contendo orientações de saúde. O material de entrada pode incluir, por exemplo, instruções sobre medicamentos, orientações gerais de acompanhamento, explicações sobre sintomas, lembretes, recomendações previamente definidas ou descrições de atividades a serem realizadas pelo paciente. 

Uma premissa da proposta é que esse conteúdo seja previamente definido ou aprovado por profissionais ou equipes de saúde. A solução não tem como objetivo verificar se o texto original está clinicamente correto, mas adaptar sua linguagem, buscando preservar o significado do material recebido. 

Além do conteúdo textual, a solução recebe informações do perfil comunicacional do paciente, como idade, escolaridade e área de formação. Esses dados são tratados como entradas para orientar a adaptação da linguagem, conforme detalhado na Seção~\ref{sec:perfil-paciente-adaptacao-comunicacional}.

A solução também considera parâmetros definidos no momento do processamento do conteúdo. Entre eles, está o nível de detalhamento desejado para a explicação. Quando a entrada é um PDF, também pode ser indicada a inclusão de imagens extraídas do próprio documento na apresentação do conteúdo. Nesse caso, as imagens não são geradas pela LLM, mas reaproveitadas a partir do material fornecido originalmente, quando forem úteis para apoiar a compreensão do paciente.

A partir dessas entradas, a solução trata a simplificação textual como uma tarefa contextualizada. Em vez de produzir uma reescrita genérica, a aplicação considera o conteúdo original, o perfil comunicacional do destinatário e os parâmetros definidos para o processamento. Com isso, busca gerar uma versão mais adequada ao paciente, mantendo a geração restrita ao material fornecido e às regras de preservação do significado original. 

\section{Perfil do paciente e adaptação comunicacional} 
\label{sec:perfil-paciente-adaptacao-comunicacional} 
O perfil do paciente é utilizado para orientar a forma como o conteúdo é apresentado, considerando que uma mesma orientação de saúde pode exigir diferentes níveis de explicação dependendo das características do usuário. Neste trabalho, o perfil comunicacional considera informações como idade, escolaridade e área de formação, que orientam escolhas de linguagem, como vocabulário, extensão das frases, explicação de termos técnicos e tom da comunicação. Assim, a solução pode utilizar frases mais curtas e vocabulário mais cotidiano para pacientes com menor familiaridade com o tema, ou manter alguns termos técnicos quando eles forem compreensíveis e relevantes ao contexto. 

Dessa forma, a adaptação proposta é linguística e comunicacional, não clínica. A solução pode reorganizar, esclarecer ou enfatizar informações presentes no texto original, mas não deve alterar seu significado nem acrescentar orientações que não estejam no conteúdo de origem.

\section{Processo de simplificação com LLMs} 
\label{sec:processo-simplificacao-llms} 

O processo de simplificação utiliza LLMs como componentes responsáveis por transformar o conteúdo original em saídas mais acessíveis ao paciente. A geração é orientada por instruções estruturadas que combinam o texto de entrada, o perfil comunicacional do paciente, os parâmetros de processamento, as regras de simplificação, os guardrails e o formato esperado para a resposta. 

A simplificação busca reduzir a complexidade linguística do texto, preservando seu significado. Entre as estratégias esperadas estão a substituição de jargões por palavras mais comuns, a explicação de termos técnicos, a redução de frases longas, a organização das informações em uma sequência mais direta e a manutenção de uma estrutura lógica clara. Quando necessário, termos técnicos importantes podem ser preservados, desde que acompanhados de explicações simples. Essa escolha está alinhada a estudos que investigam o uso de LLMs para tornar informações médicas mais acessíveis a pacientes, especialmente em contextos de educação em saúde, simplificação de conteúdos médicos e apoio ao letramento em saúde em aplicações mHealth \cite{jeblick2024chatgpt,aydin2024llms,loughran2024mhealth}.

Na solução desenvolvida, o processamento ocorre em dois fluxos de geração. O primeiro é voltado à produção do texto simplificado e adaptado ao perfil do paciente. O segundo é voltado à geração de perguntas frequentes e respostas baseadas no conteúdo original. Esses fluxos podem ser executados de forma paralela, pois partem do mesmo material de entrada, mas produzem saídas com finalidades distintas: uma versão reescrita do conteúdo e um conjunto auxiliar de perguntas e respostas para apoiar a compreensão. Além desses fluxos iniciais, a solução disponibiliza um chat contextual, no qual perguntas feitas pelo usuário são respondidas pela LLM apenas quando podem ser tratadas com base no conteúdo fornecido. 

A LLM, nesse processo, não é utilizada como mecanismo de consulta aberta nem como fonte autônoma de conhecimento médico. O modelo recebe uma tarefa delimitada: reescrever um conteúdo fornecido em linguagem mais clara e adequada ao perfil comunicacional do paciente, permanecendo restrito ao material de origem. Essa delimitação é reforçada pelos guardrails descritos na próxima subseção.

\section{Guardrails para preservação da precisão médica} 
\label{sec:guardrails-preservacao-precisao-medica} 

Os guardrails são parte central da proposta, pois estabelecem regras e restrições para orientar o comportamento da LLM durante a simplificação textual, a geração de perguntas e respostas e o funcionamento do chat contextual. Essa decisão está alinhada à ideia de que mecanismos de salvaguarda para LLMs devem ser projetados de forma sistemática, considerando o contexto de uso e os tipos de falhas que se busca mitigar \cite{dong2024guardrails}. No caso deste trabalho, sua principal função é reduzir o risco de que o modelo adicione informações não presentes no material de origem, interprete incorretamente uma orientação ou modifique o significado do conteúdo original.

Como a solução permite diferentes níveis de detalhamento, preservar a precisão médica não significa necessariamente reproduzir todas as informações com o mesmo grau de extensão do texto original. A versão gerada pode ser mais resumida ou mais detalhada, conforme o parâmetro definido no processamento. Entretanto, essa adaptação não deve contradizer o conteúdo original, criar recomendações novas, alterar orientações clínicas ou remover informações essenciais para a compreensão segura do texto. 

Os guardrails da solução incluem restrições como manter a geração vinculada ao conteúdo fornecido, não adicionar informações externas, não criar recomendações médicas novas e não modificar o sentido de instruções presentes no material de entrada. Quando o texto original contém informações como dosagens, frequências, contraindicações ou instruções de segurança, esses elementos devem ser preservados de forma compatível com o nível de detalhamento selecionado. Nas perguntas e respostas, a mesma lógica é aplicada: as respostas devem ser baseadas no conteúdo original e não devem preencher lacunas com conhecimento externo. 

Em aplicações de saúde, esse cuidado é particularmente relevante, pois a utilização de IA generativa exige atenção adicional à factualidade, à segurança e à confiabilidade das informações produzidas \cite{gangavarapu2024healthcare}. Nesse contexto, os guardrails funcionam como mecanismos de mitigação de risco, mas não como garantia absoluta de correção.

\section{Saídas da solução} 
\label{sec:saidas-solucao} 

A principal saída da solução é o texto simplificado e adaptado ao perfil comunicacional do paciente. Esse texto deve apresentar a mesma orientação de saúde do conteúdo original, mas em linguagem mais clara, direta e compatível com as características do usuário. A versão simplificada busca apoiar a compreensão do paciente sem introduzir instruções clínicas novas ou diferentes. 

Além do texto simplificado, a solução gera recursos auxiliares de compreensão. Um desses recursos é o glossário de termos técnicos, que identifica termos presentes no texto original e apresenta definições em linguagem simples. Esse recurso permite manter termos importantes quando necessário, mas acompanhados de explicações mais acessíveis. 

A solução também gera perguntas frequentes e respostas baseadas no conteúdo original. Essas perguntas têm como finalidade antecipar dúvidas comuns do paciente sobre a orientação apresentada, sem ampliar o conteúdo clínico do material de entrada. Assim, as respostas devem permanecer restritas ao texto fornecido e respeitar os guardrails definidos na solução. 

Além das perguntas frequentes geradas automaticamente, a solução oferece um recurso de chat contextual para que o usuário possa fazer perguntas sobre o conteúdo apresentado. Esse chat também é delimitado pelo material de origem: a LLM deve responder apenas quando a pergunta estiver relacionada ao texto carregado e houver informação suficiente para elaborar uma resposta com base nesse conteúdo. Caso a pergunta extrapole o escopo do material, a solução deve evitar responder com conhecimento externo e indicar que não há informações suficientes no texto fornecido. 

Por fim, a solução calcula o Flesch-PT para o texto original e para o texto simplificado. Esse indicador permite comparar as duas versões do ponto de vista da legibilidade textual, oferecendo um apoio quantitativo à análise da simplificação produzida.

A solução também disponibiliza dois recursos de acessibilidade voltados a ampliar o alcance da experiência para perfis com menor letramento digital ou dificuldade de interação por texto: síntese de fala do conteúdo simplificado, que permite ao usuário ouvir o material em vez de apenas lê-lo, e entrada por voz no chat contextual, que permite formular perguntas oralmente sem necessidade de digitação. Esses recursos não fazem parte do fluxo principal de simplificação e não foram avaliados neste trabalho, mas integram a implementação funcional como apoio à acessibilidade. 

\section{Indicador de legibilidade textual} 
\label{sec:indicador-legibilidade-textual} 

A solução inclui um indicador automático de legibilidade para apoiar a avaliação da simplificação linguística. Neste TCC, é utilizado o Índice de Legibilidade de Flesch adaptado ao português brasileiro, denominado neste trabalho como Flesch-PT. A escolha desse indicador se justifica por se tratar de uma adaptação da fórmula de Flesch ao português brasileiro \cite{martins1996readability}, conforme discutido na Seção~\ref{sec:legibilidade-textual}, que estima a facilidade de leitura a partir de características linguísticas mensuráveis como extensão das frases e número de sílabas por palavra. Os valores resultantes são interpretados segundo a escala apresentada na Tabela~\ref{tab:classificacao-flesch-pt}. 

No contexto da solução proposta, o Flesch-PT é calculado para o texto original e para a versão simplificada. Essa comparação permite observar se a adaptação produzida pela LLM resultou em uma versão com maior legibilidade estimada. Assim, o indicador funciona como um apoio quantitativo à análise da simplificação textual, permitindo comparar duas versões de um mesmo conteúdo de forma objetiva. 

Entretanto, o Flesch-PT deve ser interpretado com cautela. O índice avalia aspectos formais do texto, mas não mede diretamente a compreensão real do paciente nem garante que o conteúdo médico tenha sido preservado corretamente. Por isso, a análise da solução combina o indicador automático de legibilidade com critérios de aderência aos guardrails e preservação do significado do texto original.

\section{Delimitações da solução}
\label{sec:delimitacoes-solucao}

A solução proposta possui delimitações claras quanto ao seu escopo de atuação. Ela não realiza diagnósticos, não prescreve tratamentos, não recomenda novas condutas clínicas, não substitui profissionais da saúde e não valida clinicamente o conteúdo original. Também não utiliza o perfil do paciente para personalizar decisões médicas, mas apenas para orientar a adaptação comunicacional do texto.

Sua função é apoiar a comunicação em saúde por meio da transformação linguística de conteúdos previamente definidos ou validados. A solução atua sobre a forma de apresentação da informação, buscando torná-la mais clara e adequada ao perfil do paciente, sem modificar o significado clínico do material de origem. Dessa forma, deve ser compreendida como uma ferramenta de apoio à compreensão, e não como um sistema clínico autônomo.

No escopo deste trabalho, a Takere é tratada como ambiente de referência para aplicação da abordagem, sem integração plena realizada. Além disso, a solução foi construída e demonstrada em ambiente controlado, com uso de perfis simulados e sem dados reais de pacientes. As limitações decorrentes dessas decisões são discutidas posteriormente no Capítulo~\ref{cap:limitacoes}.